\chapter{Queues, Ring Buffers and Shared Memory}\label{ch:ll:queues-ring-buffers-and-shared-memory}

In 2011 the engineers of a retail exchange published the numbers that had made them redesign their system. Passing an
event through a three-stage pipeline built on the standard Java blocking queue took 32\,757 nanoseconds per hop on
average, and its 99.99th percentile was over four milliseconds; through the \omterm{def:ll:queues-ring-buffers-and-shared-memory:ring}{ring buffer} they built instead, 52
nanoseconds on average and under 8\,192 at the 99.99th percentile (Thompson et al., 2011). The work in each stage was
the same; the queues were the latency. A trading process is a pipeline of threads (feed handler, book, strategy, gateway,
logger), and how data moves between them decides much of its tail. This chapter builds the structures that move it: the
single-producer single-consumer ring, the broadcast ring of the ``disruptor'' pattern, the \omterm{def:ll:queues-ring-buffers-and-shared-memory:seqlock}{sequence lock} for data where
only the latest value matters, and the same structures across processes in shared memory, with what to do when a
consumer cannot keep up.

\section{The single-producer single-consumer ring}

\begin{definition}[Ring buffer, single-producer single-consumer queue]\label{def:ll:queues-ring-buffers-and-shared-memory:ring}
A \emph{ring buffer}\index{ring buffer} is a fixed array of slots used circularly: a writer's position and a reader's position
increase forever and index the array modulo its size, a power of two so that the modulo is a mask. A \emph{single-producer
single-consumer queue}\index{single-producer single-consumer queue} (SPSC) is a ring with exactly one writing thread and one
reading thread; each position is written by one thread only, so no \omterm{def:ll:lock-free-programming:cas}{compare-and-swap} is needed.
\end{definition}

The producer writes the slot at its position and then publishes the new position with a release store; the consumer loads
the producer's position with acquire, reads every slot before it, and publishes its own position the same way, handing the
slots back (chapter~11). Both operations are \omterm{def:ll:lock-free-programming:progress}{wait-free}. Lamport proved such a queue correct without locks in 1983; what
decades of practice added are the details that make it fast. The two positions live on separate \omterm{def:ll:memory-hierarchy-and-caches:line}{cache lines}, so that the
producer's writes do not invalidate the consumer's position (chapter~3). And each side keeps a private copy of the other's
position and refreshes it only when the ring looks full or empty: while the ring is neither, a message costs one \omterm{def:ll:memory-hierarchy-and-caches:line}{cache
line} transfer (the slot) instead of three.

\begin{omfigure}
\begin{tikzpicture}[font=\footnotesize]
\foreach \i/\c in {0/5,1/5,2/30,3/30,4/30,5/5,6/5,7/5} {
  \pgfmathsetmacro{\a}{90 - 45*\i}
  \draw[fill=omDef!\c,draw=omDef] (\a+22.5:1.3) arc (\a+22.5:\a-22.5:1.3) -- (\a-22.5:2.3) arc (\a-22.5:\a+22.5:2.3) -- cycle;
  \node at (\a:1.8) {\i};
}
\node[align=center] at (0,0) {8 slots\\ (mask 7)};
\draw[-Stealth,very thick,omThm] (3.9,0.9) node[right,align=left]{read = 2\\ (consumer's line)} -- (0:2.4);
\draw[-Stealth,very thick,omMeth] (-3.9,-2.3) node[left,align=right]{write = 5\\ (producer's line)} -- (-135:2.4);
\node[anchor=west,align=left] at (3.9,-1.0) {shaded: slots 2, 3, 4,\\ written, not yet read};
\node[anchor=west,align=left] at (3.9,-1.9) {full when write $-$ read $=$ 8};
\end{tikzpicture}

\omcaption{A ring of eight slots after five messages were written and two read (the state of the build's shared fixture).
The producer's and consumer's positions each sit on a \omterm{def:ll:memory-hierarchy-and-caches:line}{cache line} of their own and only grow; the slot of position $p$ is
$p \bmod 8$.}\label{fig:ll:queues-ring-buffers-and-shared-memory:ring}
\end{omfigure}

\Cref{fig:ll:queues-ring-buffers-and-shared-memory:rtt} measures a round trip between two pinned threads: one sends an
eight-byte message through a ring, the other sends it back through a second ring. About \qty{200}{\nano\second} at the
median: four cache-line transfers between cores. The same round trip through a queue protected by a mutex and a condition
variable costs about \qty{30}{\micro\second}, over a hundred times more, because the waiting thread goes to sleep in the
kernel and must be woken. One-way, the ring streams about 250 million eight-byte messages a second on this laptop.

\begin{omfigure}
\begin{tikzpicture}
\begin{axis}[width=12cm,height=5.8cm,ybar,bar width=7pt,ymode=log,log origin y=infty,ymin=100,ymax=5000000,
  symbolic x coords={threads,other-pair,processes,locked},
  xtick=data,xticklabels={{SPSC ring,\\ threads},{SPSC ring,\\ other CPU pair},{SPSC ring,\\ two processes},{mutex and\\ condition variable}},
  xticklabel style={align=center,font=\footnotesize},enlarge x limits=0.15,
  ylabel={ns per round trip (log scale)},tick label style={font=\small},label style={font=\small},grid=major,
  grid style={gray!25},legend columns=4,legend style={font=\footnotesize,at={(0.5,-0.33)},anchor=north,draw=none,fill=none,
  /tikz/every even column/.append style={column sep=6pt}},table/col sep=comma]
\addplot[fill=omDef!70,draw=omDef] table[x=key,y=p50]{figdata/low-latency/12-queues-ring-buffers-and-shared-memory/measured_rtt.csv};
\addplot[fill=omDef!35,draw=omDef] table[x=key,y=p99]{figdata/low-latency/12-queues-ring-buffers-and-shared-memory/measured_rtt.csv};
\addplot[fill=omThm!45,draw=omThm] table[x=key,y=p999]{figdata/low-latency/12-queues-ring-buffers-and-shared-memory/measured_rtt.csv};
\addplot[fill=gray!40,draw=gray] table[x=key,y=max]{figdata/low-latency/12-queues-ring-buffers-and-shared-memory/measured_rtt.csv};
\legend{p50,p99,p99.9,maximum}
\end{axis}
\end{tikzpicture}

\omcaption{Round trips of an eight-byte message between two pinned threads (or processes) through two SPSC rings, and through
two queues protected by a mutex and a condition variable; 100\,000 round trips each. Measured on a laptop (Intel Core Ultra 7
155H) under WSL2, no isolated cores. Data: \texttt{bench\_rings.py}.}\label{fig:ll:queues-ring-buffers-and-shared-memory:rtt}
\end{omfigure}

\section{Many readers: the disruptor pattern}

\begin{definition}[Disruptor pattern]\label{def:ll:queues-ring-buffers-and-shared-memory:disruptor}
In the \emph{disruptor pattern}\index{disruptor pattern} one ring carries every event to several consumers, each reading at its
own position; consumers that depend on each other wait on each other's positions rather than on queues between them, so a
pipeline of stages becomes one ring with several cursors and no copies.
\end{definition}

Market data is the natural case: one feed handler, several consumers (the book builder, a second strategy, the recorder,
a risk monitor), each needing every message. Two policies differ in who waits. With \omterm{def:ll:queues-ring-buffers-and-shared-memory:bp}{back-pressure} (next section) the
producer waits for the slowest consumer, which is right for an order path, where nothing may be lost. For market data the
producer must not wait: the exchange will not slow down. The build's broadcast ring therefore never blocks its producer; each
slot carries the sequence number of the message it holds, so a reader that has been lapped sees a later sequence than it
expected and knows it lost messages, instead of reading a mixture.

\begin{omfigure}
\begin{tikzpicture}[font=\footnotesize]
\foreach \i/\c in {0/8,1/9,2/10,3/11,4/12,5/13,6/14,7/15} {
  \draw[fill=omDef!10,draw=omDef] (1.3*\i,0) rectangle ++(1.2,0.8);
  \node at (1.3*\i+0.6,0.52) {slot \i};
  \node[font=\scriptsize] at (1.3*\i+0.6,0.2) {msg \c};
}
\draw[-Stealth,very thick,omMeth] (0.6,1.6) node[above]{producer: next = 16, into slot 0} -- (0.6,0.85);
\draw[-Stealth,thick,omDef] (7.1,-0.8) -- (7.1,-0.05);
\node[anchor=north east,align=right,omDef] at (7.0,-0.75) {recorder\\ at 13};
\draw[-Stealth,thick,omDef] (9.7,-0.8) -- (9.7,-0.05);
\node[anchor=north west,align=left,omDef] at (9.8,-0.75) {book builder\\ at 15};
\draw[-Stealth,thick,omThm] (8.4,-1.95) -- (8.4,-0.05);
\node[anchor=north,align=center,omThm] at (8.4,-1.95) {slow monitor at 6: slot 6\\ holds msg 14, lapped};
\end{tikzpicture}

\omcaption{A broadcast ring of eight slots after 16 messages. Each reader keeps its own position; the producer never waits. A
reader at position 6 finds slot 6 holding message 14, a later sequence than it expects: it has been lapped (a \omterm{def:ll:queues-ring-buffers-and-shared-memory:bp}{slow consumer})
and must resynchronise.}\label{fig:ll:queues-ring-buffers-and-shared-memory:broadcast}
\end{omfigure}

\section{Sequence locks for last-value data}

\begin{definition}[Sequence lock]\label{def:ll:queues-ring-buffers-and-shared-memory:seqlock}
A \emph{sequence lock}\index{sequence lock} publishes a small value that one writer updates and many readers read: the writer
increments a counter to an odd value, writes the value, and increments it again to an even value; a reader reads the
counter, copies the value, reads the counter again, and retries if the two readings differ or are odd.
\end{definition}

The Linux kernel documentation describes the mechanism as having lockless readers with read-only retry loops and no writer
starvation, for data rarely written relative to how often it is read. In trading the pattern suits the top of the book or a
reference price, where readers want the latest consistent value, not every intermediate one. Two details make it correct in
C++ as well as on x86: the value must be copied with \omterm{def:ll:lock-free-programming:atomic}{atomic operations} (relaxed ones suffice), since a copy that races with
the writer's non-atomic stores is a \omterm{def:ll:rust-for-low-latency:race}{data race} and \omterm{def:ll:cpp-for-latency-iii-the-compiler:ub}{undefined behaviour} (Boehm, 2012); and the reader's second counter read
must be ordered after the copy by an acquire fence. The build's test runs a writer in a tight loop against a reader that
checks two hundred thousand pairs for tearing.

\section{Shared memory between processes}

\begin{definition}[Shared-memory transport]\label{def:ll:queues-ring-buffers-and-shared-memory:shm}
A \emph{shared-memory transport}\index{shared-memory transport} carries messages between processes through a region of memory
that each maps into its address space (on Linux, a POSIX shared-memory object from \texttt{shm\_open}, or a file on a
memory filesystem), laid out as rings whose positions are atomics; no system call is made per message.
\end{definition}

Processes isolate failures: a strategy that crashes need not take the feed handler with it, and components written in
different languages can share a feed. The ring does not care: \cref{fig:ll:queues-ring-buffers-and-shared-memory:rtt} shows
the same round trip between two processes as between two threads. What changes is the contract between the processes: the
layout must be specified to the byte and versioned, since two binaries built at different times read it; every field must
be of fixed size and alignment; and nothing in the region may be a pointer, which means nothing in the other process.
The build's layout is a 192-byte header (magic number, version, slot size, capacity, then the two positions on their own
lines) followed by the slots, and its C++ and Rust implementations both reproduce the same byte image of a ring, written
independently from the specification by a script.

\section{Back-pressure, slow consumers and conflation}

\begin{definition}[Back-pressure, slow consumer, conflation]\label{def:ll:queues-ring-buffers-and-shared-memory:bp}
\emph{Back-pressure}\index{back-pressure} is the propagation of a consumer's slowness to its producer, which must wait or
refuse new work when the queue between them is full. A \emph{slow consumer}\index{slow consumer} is a consumer whose rate is
below its input's, so that its queue grows or, on a ring that does not wait, it is lapped. \emph{Conflation}\index{conflation}
is the deliberate delivery of only the latest value of each item (a price level, a quote) to a consumer that cannot take every
update.
\end{definition}

A queue's size is a statement about bursts. One Quant Book~4, chapter~8, gives the M/M/1 answer for a stage with Poisson
arrivals at rate $\lambda$ and exponential service at rate $\mu$: with utilisation $\rho = \lambda/\mu$, the stationary number
in the system exceeds $n$ with probability $\rho^n$. A ring that must overflow less than once in a billion messages at
$\rho = 0.9$ therefore needs $n$ with $0.9^n < 10^{-9}$, about 197 slots, 256 as a power of two. Market data is burstier than
Poisson: the sizing starts from the M/M/1 figure and is checked against recorded bursts (chapter~18).

\begin{proposition}[Time before a slow consumer is lapped]\label{prop:ll:queues-ring-buffers-and-shared-memory:lap}
A broadcast ring of $N$ slots whose producer writes at rate $\lambda_b$ during a burst, read by a consumer at rate $\mu < \lambda_b$
that starts the burst with an empty backlog, laps the consumer after $N/(\lambda_b - \mu)$ seconds.
\end{proposition}

\begin{proof}
The consumer's backlog grows at $\lambda_b - \mu$ messages a second; the producer overwrites the consumer's next unread slot when
the backlog reaches $N$.
\end{proof}

A ring of 65\,536 slots, a burst at 2 million messages a second and a monitor that reads 1.5 million gives the monitor 131
milliseconds; an opening burst lasting longer laps it. The design answers are to give such consumers conflated views (a
\omterm{def:ll:queues-ring-buffers-and-shared-memory:seqlock}{sequence lock} per instrument instead of every message), to size their rings for measured bursts, or to let them fall behind
by design and resynchronise from a snapshot, as the feed handler of chapter~18 does from the exchange.

\section{Tutorial: rings in threads, processes and two languages}\label{tut:ll:queues-ring-buffers-and-shared-memory:tut}

\begin{tutorial}
\textbf{Goal.} Build the rings, check the shared layout in C++ and Rust against one fixture, and measure round trips and
throughput. \textbf{End state:} \cref{fig:ll:queues-ring-buffers-and-shared-memory:rtt}, a throughput figure, and green tests in
both languages.
\begin{tutsteps}
\item \textbf{Write and publish.} The producer copies the message into its slot and publishes its position with a release store;
it reads the consumer's position only when the ring looks full.
\omcode{code/firm/ring/cpp/firm_ring.hpp}{68}{95}{The SPSC ring's producer and consumer.}\label{lst:ll:queues-ring-buffers-and-shared-memory:spsc}
\item \textbf{The same in Rust}, over the same bytes: a region of atomic words, written through their interior mutability.
\omcode{code/firm/ring/rust/src/lib.rs}{80}{101}{The Rust producer, writing the layout the C++ reader expects.}\label{lst:ll:queues-ring-buffers-and-shared-memory:rust}
\item \textbf{A \omterm{def:ll:queues-ring-buffers-and-shared-memory:seqlock}{sequence lock}} whose value is copied as relaxed atomic words.
\omcode{code/firm/ring/cpp/firm_ring.hpp}{155}{179}{Writer and reader of the sequence lock.}\label{lst:ll:queues-ring-buffers-and-shared-memory:seqlock}
\item \textbf{Check the layout.} \texttt{make\_ring\_fixture.py} writes the byte image of a ring from the specification; both
test suites build the same ring and compare, then consume the fixture.
\item \textbf{Measure} with \texttt{python bench\_rings.py}: round trips between two threads on two pairs of CPUs, between two
processes, and through a locked queue, then one-way throughput.
\end{tutsteps}
\textbf{What to change next.} Remove the cached copies of the other side's position and measure the throughput again; put the two
positions on the same \omterm{def:ll:memory-hierarchy-and-caches:line}{cache line} and measure the round trip.
\end{tutorial}

\section{Build: the firm's rings}\label{bld:ll:queues-ring-buffers-and-shared-memory:ring}

\begin{build}
\bfield{Purpose} The transport of every Part~IV component: the feed handler publishes normalised events on a broadcast ring
(chapter~18), the strategy engine reads them and hands orders to the gateway through an SPSC ring (chapters~20--21), the binary
logger drains a ring (chapter~23), and the top of each book is published with a \omterm{def:ll:queues-ring-buffers-and-shared-memory:seqlock}{sequence lock}.
\bfield{Interface} C++20 \texttt{firm::ring}: \texttt{format(base, capacity, slot\_size)}, \texttt{region\_size},
\texttt{Spsc(base)} with \texttt{try\_write(msg, len)} and \texttt{try\_read(out, cap)}; \texttt{Broadcast(base)} with
\texttt{write}, \texttt{read(pos, out, len)} returning \texttt{Ok}, \texttt{Empty} or \texttt{Overrun}, and \texttt{head()};
\texttt{SeqLock<T>} with \texttt{store} and \texttt{load}; \texttt{Segment(name, bytes, create)}. Rust \texttt{firm\_ring}:
\texttt{Region}, \texttt{format}, \texttt{Spsc::attach} with \texttt{try\_write} and \texttt{try\_read}, \texttt{SeqLock2}.
\bfield{Rules} The byte layout above, little-endian, versioned by a magic number and a version; capacities are powers of two;
positions on separate \omterm{def:ll:memory-hierarchy-and-caches:line}{cache lines}; the SPSC ring reports full and empty, the broadcast ring never waits and reports overruns;
no allocation after construction.
\bfield{Acceptance tests} \texttt{code/firm/ring/}: the fixture's bytes reproduced and consumed in C++ and Rust; 200\,000
messages in order across two threads in both; broadcast overrun detection and resynchronisation; no torn read in 200\,000
sequence-lock reads under a writer in a tight loop; 1\,000 messages from a child process over shared memory.
\bfield{Stretch} A broadcast ring with a declared slowest-consumer policy (wait, drop, conflate); a ring in a file on a huge-page
filesystem.
\end{build}

\begin{omsources}
\item M.~Thompson, D.~Farley, M.~Barker, P.~Gee and A.~Stewart, \emph{Disruptor: high performance alternative to bounded queues
for exchanging data between concurrent threads}, 2011.
\item L.~Lamport, ``Specifying concurrent program modules'', \emph{ACM TOPLAS} 5(2), 1983.
\item H.-J.~Boehm, ``Can seqlocks get along with programming language memory models?'', MSPC 2012; Linux kernel documentation,
``Sequence counters and sequential locks''.
\item Linux man page \texttt{shm\_open(3)}.
\end{omsources}

\section{Exercises}

\begin{exercise}[$\star$]\label{exo:ll:queues-ring-buffers-and-shared-memory:1}
A ring has 1\,024 slots; the producer's position is 70\,001 and the consumer's 69\,500. How many messages are waiting, which slot
will the producer write next, and how many more can it write before the ring is full?
\end{exercise}

\begin{exercise}[$\star$]\label{exo:ll:queues-ring-buffers-and-shared-memory:2}
At 250 million messages a second, how long does a 64-bit position take to wrap? And a 32-bit one?
\end{exercise}

\begin{exercise}[$\star$]\label{exo:ll:queues-ring-buffers-and-shared-memory:3}
A stage runs at $\rho = 0.8$. What ring size makes an M/M/1 overflow rarer than one in a million?
\end{exercise}

\begin{exercise}[$\star\star$]\label{exo:ll:queues-ring-buffers-and-shared-memory:4}
Why does each side of the SPSC ring keep a private copy of the other's position? What does removing the copies cost per message?
\end{exercise}

\begin{exercise}[$\star\star$]\label{exo:ll:queues-ring-buffers-and-shared-memory:5}
A risk monitor reads the broadcast ring at 400\,000 messages a second. A burst of 3 million messages a second lasts 80
milliseconds on a ring of 131\,072 slots. Is the monitor lapped, and if so when?
\end{exercise}

\begin{exercise}[$\star\star$]\label{exo:ll:queues-ring-buffers-and-shared-memory:6}
Why must the \omterm{def:ll:queues-ring-buffers-and-shared-memory:seqlock}{sequence lock}'s value be copied with \omterm{def:ll:lock-free-programming:atomic}{atomic operations} in C++, although the reader discards any copy taken during a
write?
\end{exercise}

\begin{exercise}[$\star\star\star$]\label{exo:ll:queues-ring-buffers-and-shared-memory:7}
\emph{Coding.} Put the producer's and consumer's positions on the same \omterm{def:ll:memory-hierarchy-and-caches:line}{cache line} (edit the header offsets in a copy of
\texttt{firm\_ring.hpp}) and rerun the round-trip and throughput benchmarks. What changes, and why?
\end{exercise}

\begin{exercise}[$\star\star\star$]\label{exo:ll:queues-ring-buffers-and-shared-memory:8}
\emph{Find the flaw.} ``Our shared-memory ring stores a pointer to the head message and a \texttt{std::string} for the instrument
name in its header.''
\end{exercise}

\section{Problem: Three Strategies, One Feed}

\begin{problem}[{Weekend problem --- sizing a broadcast ring}]\label{pb:ll:queues-ring-buffers-and-shared-memory:1}
A feed handler publishes normalised events on a broadcast ring of 65\,536 slots of 64 bytes. Three consumers read it: a book
builder at up to 4 million events a second, a strategy at 2.5 million, and a monitor at 1.5 million. On an ordinary day the feed
averages 300\,000 events a second; at the open, bursts reach 2 million a second for up to half a second.

\medskip\noindent\textbf{Part I --- The ring.}
\begin{enumerate}
\item How many bytes does the ring occupy, and does it fit the second-level cache of chapter~3?
\item What is each consumer's utilisation on an ordinary day?
\item Which consumers can fall behind during a burst?
\item Why does the producer never wait for them?
\end{enumerate}

\medskip\noindent\textbf{Part II --- The burst.}
\begin{enumerate}[resume]
\item How long before the monitor is lapped (\cref{prop:ll:queues-ring-buffers-and-shared-memory:lap})?
\item Is it lapped during a half-second burst?
\item What ring size would carry the monitor through the burst?
\item How much memory would that cost?
\end{enumerate}

\medskip\noindent\textbf{Part III --- The alternatives.}
\begin{enumerate}[resume]
\item The monitor needs only the latest price of each of 5\,000 instruments. What does a \omterm{def:ll:queues-ring-buffers-and-shared-memory:seqlock}{sequence lock} per instrument cost in memory,
and what share of updates does the monitor skip during the burst?
\item What does the monitor do when it detects an overrun on the ring?
\item Why is \omterm{def:ll:queues-ring-buffers-and-shared-memory:bp}{back-pressure} wrong for this producer and right for the order path?
\item What does the M/M/1 model say about the strategy's ring at 300\,000 events a second, and why is it optimistic?
\end{enumerate}

\medskip\noindent\textbf{Part IV --- The verdict.}
\begin{enumerate}[resume]
\item State the \emph{named result}: the time before the monitor is lapped in a burst, and the share of updates a conflating monitor
skips.
\item What does the measured round trip say about the cost of adding a consumer?
\item Why is the ring's layout versioned?
\item What would the published comparison of a blocking queue and a ring predict for a blocking queue here?
\item Which placement rule of chapter~4 applies to the feed handler and the book builder?
\item How would you record the burst sizes that the ring must carry?
\item What must the strategy do after a gap in its input?
\item In one sentence: what is a \omterm{def:ll:queues-ring-buffers-and-shared-memory:ring}{ring buffer}'s size a statement about?
\end{enumerate}
\end{problem}

\section{Interview questions}

\begin{interviewq}[$\star$ \iqroles{developer}]\label{iq:ll:queues-ring-buffers-and-shared-memory:1}
Implement a single-producer single-consumer \omterm{def:ll:queues-ring-buffers-and-shared-memory:ring}{ring buffer}. Which \omterm{def:ll:lock-free-programming:ordering}{memory orderings} does it need?
\end{interviewq}

\begin{interviewq}[$\star\star$ \iqroles{developer}]\label{iq:ll:queues-ring-buffers-and-shared-memory:2}
Why is a \omterm{def:ll:queues-ring-buffers-and-shared-memory:ring}{ring buffer} faster than a queue protected by a mutex and a condition variable?
\end{interviewq}

\begin{interviewq}[$\star\star$ \iqroles{developer}]\label{iq:ll:queues-ring-buffers-and-shared-memory:3}
What is a \omterm{def:ll:queues-ring-buffers-and-shared-memory:seqlock}{sequence lock}, when would you use one, and what are its pitfalls?
\end{interviewq}

\begin{interviewq}[$\star\star$ \iqroles{developer}]\label{iq:ll:queues-ring-buffers-and-shared-memory:4}
One market-data feed, three consumers of different speeds. Design the distribution.
\end{interviewq}

\begin{interviewq}[$\star\star$ \iqroles{developer}]\label{iq:ll:queues-ring-buffers-and-shared-memory:5}
What must a data structure in shared memory between two processes avoid?
\end{interviewq}

\begin{interviewq}[$\star\star\star$ \iqroles{developer, researcher}]\label{iq:ll:queues-ring-buffers-and-shared-memory:6}
How would you size the queue in front of a stage whose input is bursty? What does queueing theory give you, and what does it
miss?
\end{interviewq}
